% English translation of Hieronymus Georg Zeuthen's "Note sur la théorie de surfaces réciproques" (1871).
% Translated from the transcription of the original French text (original.tex).
% All formulas, tags, markers, and math structures are preserved verbatim.
\documentclass{article}
\usepackage{readmasters}
\begin{document}

\origpage{633}
\section*{Note on the Theory of Reciprocal Surfaces.}

\subsection*{By H.~G. Zeuthen in Copenhagen.}

In my two memoirs at the beginning of this volume of the „Annalen“, I attributed to the surfaces under consideration certain singularities that I enumerated in the introduction to the first of these memoirs. I assumed that no other singularity was added, or else that the singularities presented themselves in the most general manner. Now, as this most general manner differs according to whether one regards the surface as a locus of points or as an envelope of planes, one must divide the singularities into \emph{singularities of punctual definition} and \emph{singularities of tangential definition}. Singularities of punctual definition are those which satisfy their definition in the most general manner if one regards the surface as a locus of points, and singularities of tangential definition are those which satisfy it in the most general manner if one regards the surface as an envelope of planes.

We did not state explicitly which definitions of singularities are punctual and which are tangential, without which the definitions are not complete. We shall supply this omission here by showing that the division we followed, without stating it, is the most natural one. Since a surface regarded as a locus of points does not, in general, possess double and cuspidal curves, but rather series of double and stationary tangent planes, the definitions of the ordinary singularities of the double and stationary tangent planes ($b', k', t', \gamma', q', \varrho', c', h', \beta', r', \sigma'$) must be punctual, and those of the ordinary singularities of the double curve and of the cuspidal curve ($b, k, t, \gamma, q, \varrho, c, h, \beta, r, \sigma$) must be tangential.${}^{*)}$ A

\textbf{*)} In that case one knows that the $\beta$ cuspidal points of the cuspidal curve lie on the double curve, and that the $\gamma$ cuspidal points of the double curve lie on the cuspidal curve, and these points are therefore the same as those whose numbers Messrs.~Salmon and Cayley designate in the same manner. Mr.~Cayley has observed to me that the two curves may have other cuspidal points. These new singular points would require a new study. --- I ought to have stated explicitly that my definition of $k$ (the Plückerian number of double generators of a cone projecting the double curve) differs somewhat from that of Messrs.~Salmon and Cayley, which explains an apparent discrepancy between their formulas and mine. I write $k - f - 3t - \Sigma(\cdots)$ where Mr.~Cayley writes simply $k$.

\origpage{634}
double curve, once attributed to a surface regarded as a locus of points, will in general be endowed with pinch-points, whereas the double curve of a surface regarded as an envelope of planes generally has none. For this reason one must give these $j$ singularities a punctual definition, and the $j'$ pinch-planes a tangential definition. Likewise the $\chi$ close-points are of punctual definition, and the $\chi'$ close-planes of tangential definition. We explicitly gave a punctual definition to the conical points and to the singularities belonging to them, and a tangential definition to the planes that are tangent along a curve and to the singularities belonging to them. It is by regarding the definitions of $f$ and $i$ as punctual, and those of $f'$ and $i'$ as tangential, that one obtains the equations $f = f'$ and $i = i'$.

We make yet another distinction: the punctual properties of a singularity are those relating to the points of the surface, and the tangential properties are those relating to its tangent planes. For example, $t'$ denotes the number of planes having the punctual property of meeting the surface in a curve endowed with three double points, and the tangential property of being triple tangent planes, either of the surface itself or of the envelope of the double tangent planes. The punctual properties of singularities of punctual definition, as well as the tangential properties of singularities of tangential definition, are indicated in the definitions themselves; but it is often very difficult to find the punctual properties of singularities of tangential definition, and the tangential properties of singularities of punctual definition, and these properties can be quite complicated. In order to know sufficiently the tangential properties of a pinch-point, for example, one must have studied the singularity of the envelope of stationary tangent planes caused by the pinch-point. To find it, one needs to study the singularities that the Hessian and its curve of intersection with the given surface have at that same point.

If, in an investigation, one regards a surface as the locus of its points, one must know the punctual properties of the singularities attributed to it, and if one regards a surface as the envelope of its tangent planes, one must know the tangential properties of its singularities. In a theory of reciprocal surfaces, one needs to know, at once, both the punctual properties and the tangential properties of the singularities attributed to the surfaces. Having made a sufficiently deep study neither of the tangential properties of the singularities $j, \chi, \eta$ and $\xi$, whose definition is punctual, nor of the punctual properties of the singularities $j', \chi', \eta'$ and $\xi'$, whose definition is tangential, I must request the readers of my memoirs \emph{to append to the}

\origpage{635}
memoir „Sur deux surfaces dont les points se correspondent un-à-un“, where the surfaces are regarded as loci of points, the following restriction:

„\emph{The surfaces $[F_1]$ and $[F_2]$ are devoid of the singularities $j', \chi', \eta'$ and $\xi'$“},

and to the introduction of the memoir „Sur les droites multiples des surfaces“ the following restriction:

„\emph{The equations on page 4\ednote{In the printed text the digit 4 has a heavy ink defect or batter; the context of the memoir establishes page 4.} are correct if the surface is devoid of the singularities\ednote{Printed singularites without accent; an apparent misprint for singularités.} $j', \chi', \eta'$ and $\xi'$, and the reciprocal equations if it is devoid of the singularities $j, \chi, \eta$ and $\xi$. They will therefore be applicable at once to a surface devoid of all these singularities.}\ednote{The quotation opened with „Les équations lacks a closing quotation mark in the print.}

At neither of these places will the restrictions alter the statements of the results; but it is possible that the singularities of which I speak, if the surfaces possessed them, would introduce new terms into the equations. As for the investigations which form the principal subject of my first memoir, I have no need to add a restriction there as well; for the singularities $j$ and $j'$ present themselves there only in a special manner.

I do not say that the formulas on page 4 cease to be true beyond the limits indicated here, but only that \emph{my deduction} of these formulas (which I did not set forth, but in which I make use either of the procedures of Mr.~Salmon or of the principle of correspondence) is then insufficient. For the singularities $j$ and $j'$, $\chi$ and $\chi'$, I merely reproduce the formulas of Mr.~Cayley (On reciprocal Surfaces. Phil.\ Trans.\ 1869, p.~202), and as for the singularities $\eta$ and $\eta'$, my results agree with those which that scholar found in relation to „binodes“ and „biplanes“; but I do not know whether Mr.~Cayley intends to give his definitions the same degree of generality that I have indicated here, or whether he restricts himself to cases where singularities of tangential definition have no punctual properties other than those which already belong to the punctual singularities he attributes to the surfaces, and conversely. This holds for cubic surfaces, which form a principal subject of his researches; but already in the study of a surface of the $4^{\text{th}}$ order endowed with a double conic and to which one gives a punctual definition, one meets with pinch-points that have more complicated tangential properties. The numbers of singularities of this surface, derived from the equations of which I have just spoken (Cayley p.~202, Zeuthen p.~4), do not satisfy the equation given by Mr.~Cayley on page 228. We shall show in what follows another fact that will be at variance with this latter equation (the others being assumed correct) if one attributes to the surface close-points defined as I have done here.

\origpage{636}
By restricting the study of reciprocal surfaces to surfaces devoid of the singularities $j, j', \chi, \chi', \eta, \eta', \xi, \xi'$, it will be less difficult to determine all the coefficients of the $26^{\text{th}}$ equation${}^{*)}$ which, alongside the 3 equations (I), the 11 equations (II)---(IV) and (VIII)---(IX), and the 11 reciprocal equations, hold between the singularities of a surface. Indeed, this $26^{\text{th}}$ equation will be equation (13) [or (14)] on page 46 of my second memoir.

I shall give here another deduction of the same equation in which I shall not make use of the theory of two surfaces whose points correspond one-to-one, but only of the principle of correspondence. We shall apply this principle to \emph{finding the number of planes that are at once tangent to the surface and osculating the cuspidal curve at one and the same point of the latter}. We then take as corresponding points ($x$ and $u$) whose coincidences are sought the points where the tangent plane and the osculating plane at the same point of the cuspidal curve meet a fixed right line. To a point $x$ there correspond $\sigma$ points $u$, and, if $m$ denotes the class of the developable (of order $r$) envelope of the osculating planes of the cuspidal curve, to a point $u$ there correspond $m$ points $x$. There are therefore $\sigma + m$ coincidences of a point $x$ with a corresponding point $u$. These coincidences occur $1^0$ at the $r$ points where the fixed right line meets tangents to the cuspidal curve, and $2^0$ at the points where it meets the sought-for planes. The number of these planes is therefore $\sigma + m - r$.

A part of the planes found, say provisionally $\Sigma'$, coincide with the planes that are tangent along a curve; the points of contact of the others will have the property reciprocal to that of the planes: namely, that of being at once points of contact of stationary tangent planes with the surface and points of contact of the same planes with the edge of regression of the envelope of the stationary tangent planes. By determining the number of these points in an analogous manner, one finds the following equation:
\[
\sigma + m - r - \Sigma' = \sigma' + m' - r' - \Sigma,
\]
where $m'$ is the order of the edge of regression of the envelope of the stationary tangent planes, and $\Sigma$ is the number of points having the indicated property that lie at the conical points of the surface. Now by Plücker's relations one has
\[
m = \beta + 3r - 3c \quad , \quad m' = \beta' + 3r' - 3c',
\]
and $\Sigma$ will be equal to the sum of the respective numbers of right lines that have

\textbf{*)} After having set forth the other 25 equations on pages 3 and 4, I speak on page 4 (and on page 46) of \emph{two} new equations (serving to determine $\beta$ and $\beta'$); but these equations reduce, by means of the 25 others, to a single one. (See the memoir of Mr.~Cayley, p.~228.)

\origpage{637}
at the conical ($\mu$-tuple) points a $(\mu + 2)$-point contact with the surface (see on page 39), or rather to
\[
\Sigma [\mu(\mu + 1) - 2y - 3z] = \Sigma [2\mu + v].
\]
One will likewise have to substitute for $\Sigma'$ the number $\Sigma' [2\mu' + v']$. One thus finds
\[
\sigma + 2r - 3c + \Sigma(2\mu + v) = \sigma' + 2r' - 3c' + \Sigma'(2\mu' + v'),
\]
which is a form of the sought-for equation. By means of the equations on page 4, one can reduce it to equation (13) on page 46.

One could apply the same procedure to a surface endowed with the singularities $j, j'$, etc.; but not knowing sufficiently the properties of these singularities, as I have said, I could not then determine the coefficients. Nevertheless, this investigation may serve to confirm that the introduction of these singularities into the theory of reciprocal surfaces was premature. Suppose, for example, that the surface is endowed with $\chi$ close-points and $\chi'$ close-planes. Then one must replace the formula found by the following:
\[
\begin{aligned}
\sigma + 2r - 3c + \Sigma(2\mu + v) + A \cdot \chi \\
= \sigma' + 2r' - 3c' + \Sigma'(2\mu' + v') + A \cdot \chi',
\end{aligned}
\]
where $A$ is an unknown but \emph{integer} coefficient; but Mr.~Cayley's equation on page 228, combined with his other formulas (on page 202, my formulas of page 4), would assign to $\chi$ and $\chi'$ the coefficient $\frac{25}{3}$.${}^{*)}$

\textbf{*)} I take the opportunity to reply to a remark of Mr.~Nöther (Nachrichten der Königl.\ Gesellschaft der Wissenschaften zu Göttingen 1871, p.~277): In the remarks I made on page 324 of the $3^{\text{rd}}$ vol.\ of the „Mathematische Annalen“ on the representation of a surface on a double plane, I make, without stating it I admit, the same assumption that I make explicitly in my memoir on „deux surfaces dont les points se correspondent un-à-un“, namely, that none of the fundamental points is endowed with a fundamental direction.
\end{document}
